% questions.tex —— 2012 年·福建卷·文科 · 题目内容片段（无答案）
% 由 tools/generate_projects.py 依来源总账自 2012/2012_数学真题_福建文_试题.md 生成，勿手改（改源 MD 后重跑）


\begin{gknotice}
\textbf{来源说明}：本题面依据固定源 PDF 的两页 A3 题面逐题转录，原卷印刷页码为 731—732。第 731 面为选择题 1—12、填空题 13—16 及解答题 17—18；第 732 面为解答题 19—22。

\textbf{图形说明}：第 6 题程序框图、第 16 题图 1 至图 3、第 19 题长方体、第 21 题抛物线与等边三角形，均由项目内 TikZ 矢量源重绘；其余各题原页无图。

\textbf{解析归属}：固定源不附答案或解析；配套答案与解析均由本库按本题面独立复算。
\end{gknotice}

\gkpart{一、选择题}

\begin{question}
\gkti{1}复数 $(2+\mathrm i)^2$ 等于\mbox{（　　）}

（A）$3+4\mathrm i$；（B）$5+4\mathrm i$；（C）$3+2\mathrm i$；（D）$5+2\mathrm i$。

\end{question}
\begin{question}
\gkti{2}已知集合 $M=\{1,2,3,4\}$，$N=\{-2,2\}$，下列结论成立的是\mbox{（　　）}

（A）$N\subseteq M$；（B）$M\cup N=M$；（C）$M\cap N=N$；（D）$M\cap N=\{2\}$。

\end{question}
\begin{question}
\gkti{3}已知向量 $\boldsymbol a=(x-1,2)$，$\boldsymbol b=(2,1)$，则 $\boldsymbol a\perp\boldsymbol b$ 的充要条件是\mbox{（　　）}

（A）$x=-\dfrac12$；（B）$x=-1$；（C）$x=5$；（D）$x=0$。

\end{question}
\begin{question}
\gkti{4}一个几何体的三视图形状都相同，大小均相等，那么这个几何体不可以是\mbox{（　　）}

（A）球；（B）三棱锥；（C）正方体；（D）圆柱。

\end{question}
\begin{question}
\gkti{5}已知双曲线 $\dfrac{x^2}{a^2}-\dfrac{y^2}{5}=1$ 的右焦点为 $(3,0)$，则该双曲线的离心率等于\mbox{（　　）}

（A）$\dfrac{3\sqrt{14}}{14}$；（B）$\dfrac{3\sqrt2}{4}$；（C）$\dfrac32$；（D）$\dfrac43$。

\Needspace{0.4\textheight}
\end{question}
\begin{question}
\gkti{6}阅读如图所示的程序框图，运行相应的程序，输出的 $s$ 值等于\mbox{（　　）}
\nopagebreak[4]
\begin{center}% ============================================================
% f2012_fujian_arts_q06_flowchart.tex
% 2012 福建卷（文科）第 6 题 程序框图（TikZ 矢量重绘）
% 开始 → k=1,s=1 → 判断 k<4
%   是：s=2s-k → k=k+1 → 回到判断
%   否：输出 s → 结束
% 来源 PDF：普通高考/2012/2012福建文.pdf，印面 731 第 1 栏。
% ============================================================
\begin{tikzpicture}[
  x=1cm, y=1cm, line width=0.6pt, font=\footnotesize, >=Stealth,
  every node/.style={inner sep=3pt},
  proc/.style={rectangle, draw, minimum width=2.7cm, minimum height=0.6cm},
  io/.style={trapezium, draw, trapezium left angle=70, trapezium right angle=110,
             minimum width=2.5cm, minimum height=0.62cm, trapezium stretches=true},
  dec/.style={diamond, draw, aspect=2.3, minimum width=2.2cm, minimum height=0.95cm},
  term/.style={rectangle, rounded corners=8pt, draw, minimum width=1.7cm, minimum height=0.56cm}
]

  \node[term] (start) at (0,0) {开始};
  \node[proc] (init)  at (0,-1.25) {$k=1$，$s=1$};
  \node[dec]  (judge) at (0,-2.95) {$k<4$?};

  \node[proc] (sbox)  at (3.05,-2.35) {$s=2s-k$};
  \node[proc] (kbox)  at (3.05,-1.45) {$k=k+1$};

  \node[io]   (out)   at (0,-4.55) {输出 $s$};
  \node[term] (end)   at (0,-5.65) {结束};

  \draw[->] (start) -- (init);
  \draw[->] (init) -- (judge);
  \draw[->] (judge.east) -- (3.05,-2.95) -- (sbox.south);
  \draw[->] (sbox.north) -- (kbox.south);
  \draw[->] (kbox.north) -- (3.05,-0.92) -- (0,-0.92);
  \draw[->] (judge.south) -- (out.north);
  \draw[->] (out) -- (end);

  \node at (1.62,-2.74) {是};
  \node at (0.38,-3.72) {否};

\end{tikzpicture}
\end{center}
\nopagebreak[4]
（A）$-3$；（B）$-10$；（C）$0$；（D）$-2$。

\end{question}
\begin{question}
\gkti{7}直线 $x+\sqrt3y-2=0$ 与圆 $x^2+y^2=4$ 相交于 $A$，$B$ 两点，则弦 $AB$ 的长度等于\mbox{（　　）}

（A）$2\sqrt5$；（B）$2\sqrt3$；（C）$\sqrt3$；（D）$1$。

\end{question}
\begin{question}
\gkti{8}函数 $f(x)=\sin\left(x-\dfrac\pi4\right)$ 的图象的一条对称轴是\mbox{（　　）}

（A）$x=\dfrac\pi4$；（B）$x=\dfrac\pi2$；（C）$x=-\dfrac\pi4$；（D）$x=-\dfrac\pi2$。

\end{question}
\begin{question}
\gkti{9}设
\nopagebreak[4]
\[f(x)=\begin{cases}1,&x>0,\\0,&x=0,\\-1,&x<0,\end{cases}\qquad g(x)=\begin{cases}1,&x\text{ \text{为有理数}},\\0,&x\text{ \text{为无理数}},\end{cases}\]
则 $f(g(\pi))$ 的值为\mbox{（　　）}

（A）$1$；（B）$0$；（C）$-1$；（D）$\pi$。

\end{question}
\begin{question}
\gkti{10}若直线 $y=2x$ 上存在点 $(x,y)$ 满足约束条件
\nopagebreak[4]
\[\begin{cases}x+y-3\leqslant0,\\x-2y-3\leqslant0,\\x\geqslant m,\end{cases}\]
则实数 $m$ 的最大值为\mbox{（　　）}

（A）$-1$；（B）$1$；（C）$\dfrac32$；（D）$2$。

\end{question}
\begin{question}
\gkti{11}数列 $\{a_n\}$ 的通项公式 $a_n=n\cos\dfrac{n\pi}{2}$，其前 $n$ 项和为 $S_n$，则 $S_{2012}$ 等于\mbox{（　　）}

（A）$1006$；（B）$2012$；（C）$503$；（D）$0$。

\end{question}
\begin{question}
\gkti{12}已知 $f(x)=x^3-6x^2+9x-abc$，$a<b<c$，且 $f(a)=f(b)=f(c)=0$。现给出如下结论：① $f(0)f(1)>0$；② $f(0)f(1)<0$；③ $f(0)f(3)>0$；④ $f(0)f(3)<0$。其中正确结论的序号是\mbox{（　　）}

（A）①③；（B）①④；（C）②③；（D）②④。

\end{question}
\gkpart{二、填空题}

\begin{question}
\gkti{13}在 $\triangle ABC$ 中，已知 $\angle BAC=60^\circ$，$\angle ABC=45^\circ$，$BC=\sqrt3$，则 $AC=$\gkblank{}。

\end{question}
\begin{question}
\gkti{14}一支田径队有男女运动员 $98$ 人，其中男运动员有 $56$ 人。按男女比例用分层抽样的方法，从全体运动员中抽出一个容量为 $28$ 的样本，那么应抽取女运动员人数是\gkblank{}。

\end{question}
\begin{question}
\gkti{15}已知关于 $x$ 的不等式 $x^2-ax+2a>0$ 在 $\mathbf R$ 上恒成立，则实数 $a$ 的取值范围是\gkblank{}。

\Needspace{0.4\textheight}
\end{question}
\begin{question}
\gkti{16}某地区规划道路建设，考虑道路铺设方案。方案设计图中，点表示城市，两点之间连线表示两城市间可铺设道路，连线上数据表示两城市间铺设道路的费用，要求从任一城市都能到达其余各城市，并且铺设道路的总费用最小。例如：在三个城市道路设计中，若城市间可铺设道路的线路图如图 1，则最优设计方案如图 2，此时铺设道路的最小总费用为 $10$。现给出该地区可铺设道路的线路图如图 3，则铺设道路的最小总费用为\gkblank{}。
\nopagebreak[4]
\begin{center}% ============================================================
% f2012_fujian_arts_q16_fig1.tex
% 2012 福建卷（文科）第 16 题 图 1：三城市线路图
% A、B、C 三城市，AB=4，AC=6，BC=8（数据为道路铺设费用）
% 来源 PDF：普通高考/2012/2012福建文.pdf，印面 731 第 2 栏下部。
% ============================================================
\begin{tikzpicture}[>=Stealth, line width=0.7pt, font=\small, x=1cm, y=1cm]
  \coordinate (A) at (1.45,1.60);
  \coordinate (B) at (0.05,1.05);
  \coordinate (C) at (3.15,0.85);

  \draw (A) -- (B);
  \draw (A) -- (C);
  \draw (B) -- (C);

  \fill (A) circle (1.4pt) node[above] {$A$};
  \fill (B) circle (1.4pt) node[below left] {$B$};
  \fill (C) circle (1.4pt) node[right] {$C$};

  \node at (0.62,1.50) {$4$};
  \node at (2.40,1.42) {$6$};
  \node at (1.55,0.78) {$8$};
  \node at (1.55,-0.55) {图 1};
\end{tikzpicture}
\end{center}
\nopagebreak[4]
\nopagebreak[4]
\begin{center}% ============================================================
% f2012_fujian_arts_q16_fig2.tex
% 2012 福建卷（文科）第 16 题 图 2：图 1 的最优设计方案
% 连接 AB（费用 4）与 AC（费用 6），总费用 10
% 来源 PDF：普通高考/2012/2012福建文.pdf，印面 731 第 2 栏下部。
% ============================================================
\begin{tikzpicture}[>=Stealth, line width=0.7pt, font=\small, x=1cm, y=1cm]
  \coordinate (A) at (1.55,1.65);
  \coordinate (B) at (0.10,0.85);
  \coordinate (C) at (3.40,1.00);

  \draw (A) -- (B);
  \draw (A) -- (C);

  \fill (A) circle (1.4pt) node[above] {$A$};
  \fill (B) circle (1.4pt) node[below left] {$B$};
  \fill (C) circle (1.4pt) node[right] {$C$};

  \node at (0.62,1.40) {$4$};
  \node at (2.55,1.45) {$6$};
  \node at (1.60,-0.35) {图 2};
\end{tikzpicture}
\end{center}
\nopagebreak[4]
\nopagebreak[4]
\begin{center}% ============================================================
% f2012_fujian_arts_q16_fig3.tex
% 2012 福建卷（文科）第 16 题 图 3：七城市道路线路图
% 边权：AB=9，AF=3，AE=2，BF=7，BC=5，FG=1，FE=3，
%       GC=3，GD=2，GE=6，CD=6，DE=6
% 来源 PDF：普通高考/2012/2012福建文.pdf，印面 731 第 2 栏下部。
% ============================================================
\begin{tikzpicture}[>=Stealth, line width=0.7pt, font=\small, x=1cm, y=1cm]
  \coordinate (A) at (0.00,2.55);
  \coordinate (B) at (1.85,3.35);
  \coordinate (C) at (4.35,2.35);
  \coordinate (D) at (3.65,0.55);
  \coordinate (E) at (0.65,0.00);
  \coordinate (F) at (1.55,1.85);
  \coordinate (G) at (2.75,1.55);

  \draw (A) -- (B);
  \draw (A) -- (F);
  \draw (A) -- (E);
  \draw (B) -- (F);
  \draw (B) -- (C);
  \draw (F) -- (G);
  \draw (F) -- (E);
  \draw (G) -- (C);
  \draw (G) -- (D);
  \draw (G) -- (E);
  \draw (C) -- (D);
  \draw (D) -- (E);

  \fill (A) circle (1.3pt) node[left] {$A$};
  \fill (B) circle (1.3pt) node[above] {$B$};
  \fill (C) circle (1.3pt) node[right] {$C$};
  \fill (D) circle (1.3pt) node[right] {$D$};
  \fill (E) circle (1.3pt) node[below left] {$E$};
  \fill (F) circle (1.3pt) node[below left] {$F$};
  \fill (G) circle (1.3pt) node[right] {$G$};

  \node at (0.88,3.16) {$9$};
  \node at (0.55,2.18) {$3$};
  \node at (0.18,1.25) {$2$};
  \node at (1.90,2.55) {$7$};
  \node at (3.30,3.02) {$5$};
  \node at (2.16,1.84) {$1$};
  \node at (0.92,1.02) {$3$};
  \node at (3.50,2.14) {$3$};
  \node at (3.42,1.02) {$2$};
  \node at (1.58,0.64) {$6$};
  \node at (4.18,1.42) {$6$};
  \node at (2.18,-0.02) {$6$};

  \node at (2.20,-0.75) {图 3};
\end{tikzpicture}
\end{center}
\nopagebreak[4]
\end{question}
\gkpart{三、解答题}

\begin{question}
\gkti{17}在等差数列 $\{a_n\}$ 和等比数列 $\{b_n\}$ 中，$a_1=b_1=1$，$b_4=8$，$\{a_n\}$ 的前 $10$ 项和 $S_{10}=55$。

（1） 求 $a_n$ 和 $b_n$；

（2） 现分别从 $\{a_n\}$ 和 $\{b_n\}$ 的前 $3$ 项中各随机抽取一项，写出相应的基本事件，并求这两项的值相等的概率。

\end{question}
\begin{question}
\gkti{18}某工厂为了对新研发的一种产品进行合理定价，将该产品按事先拟定的价格进行试销，得到如下数据：

\begin{center}
\begin{tabular}{lllllll}
\toprule
单价 $x$（元） & $8$ & $8.2$ & $8.4$ & $8.6$ & $8.8$ & $9$ \\
\midrule
销量 $y$（件） & $90$ & $84$ & $83$ & $80$ & $75$ & $68$ \\
\bottomrule
\end{tabular}
\end{center}

（1） 求回归直线方程 $\hat y=bx+a$，其中 $b=-20$，$a=\bar y-b\bar x$；

（2） 预计在今后的销售中，销量与单价仍然服从（1）中的关系，且该产品的成本是 $4$ 元/件，为使工厂获得最大利润，该产品的单价应定为多少元？（利润 $=$ 销售收入 $-$ 成本）

\Needspace{0.4\textheight}
\end{question}
\begin{question}
\gkti{19}如图，在长方体 $ABCD-A_1B_1C_1D_1$ 中，$AB=AD=1$，$AA_1=2$，$M$ 为棱 $DD_1$ 上的一点。

（1） 求三棱锥 $A-MCC_1$ 的体积；

（2） 当 $A_1M+MC$ 取得最小值时，求证：$B_1M\perp$ 平面 $MAC$。
\nopagebreak[4]
\begin{center}% ============================================================
% f2012_fujian_arts_q19_cuboid.tex
% 2012 福建卷（文科）第 19 题 长方体 ABCD-A1B1C1D1
% AB=AD=1，AA1=2，M 为棱 DD1 上一点
% 来源 PDF：普通高考/2012/2012福建文.pdf，印面 732 第 1 栏上部。
% ============================================================
\begin{tikzpicture}[>=Stealth, line width=0.7pt, font=\small, x=1cm, y=1cm]
  \coordinate (B) at (0.00,0.00);
  \coordinate (C) at (1.60,0.00);
  \coordinate (B1) at (0.00,2.55);
  \coordinate (C1) at (1.60,2.55);
  \coordinate (A) at (0.95,0.58);
  \coordinate (D) at (2.55,0.58);
  \coordinate (A1) at (0.95,3.13);
  \coordinate (D1) at (2.55,3.13);
  \coordinate (M) at (2.55,1.55);

  % 可见棱
  \draw (B) -- (C) -- (C1) -- (B1) -- cycle;
  \draw (B1) -- (A1) -- (D1) -- (C1);
  \draw (C) -- (D) -- (D1);

  % 不可见棱
  \draw[dashed] (B) -- (A) -- (D);
  \draw[dashed] (A) -- (A1);

  \fill (A1) circle (1.2pt) node[above left] {$A_1$};
  \fill (B1) circle (1.2pt) node[left] {$B_1$};
  \fill (C1) circle (1.2pt) node[right] {$C_1$};
  \fill (D1) circle (1.2pt) node[right] {$D_1$};
  \fill (A) circle (1.2pt) node[below left] {$A$};
  \fill (B) circle (1.2pt) node[below left] {$B$};
  \fill (C) circle (1.2pt) node[below right] {$C$};
  \fill (D) circle (1.2pt) node[right] {$D$};
  \fill (M) circle (1.2pt) node[right] {$M$};
\end{tikzpicture}
\end{center}
\nopagebreak[4]
\end{question}
\begin{question}
\gkti{20}某同学在一次研究性学习中发现，以下五个式子的值都等于同一个常数：

① $\sin^2 13^\circ+\cos^2 17^\circ-\sin 13^\circ\cos 17^\circ$；

② $\sin^2 15^\circ+\cos^2 15^\circ-\sin 15^\circ\cos 15^\circ$；

③ $\sin^2 18^\circ+\cos^2 12^\circ-\sin 18^\circ\cos 12^\circ$；

④ $\sin^2(-18^\circ)+\cos^2 48^\circ-\sin(-18^\circ)\cos 48^\circ$；

⑤ $\sin^2(-25^\circ)+\cos^2 55^\circ-\sin(-25^\circ)\cos 55^\circ$。

（1） 试从上述五个式子中选择一个，求出这个常数；

（2） 根据（1）的计算结果，将该同学的发现推广为三角恒等式，并证明你的结论。

\Needspace{0.4\textheight}
\end{question}
\begin{question}
\gkti{21}如图，等边三角形 $OAB$ 的边长为 $8\sqrt3$，且其三个顶点均在抛物线 $E:x^2=2py$（$p>0$）上。

（1） 求抛物线 $E$ 的方程；

（2） 设动直线 $l$ 与抛物线 $E$ 相切于点 $P$，与直线 $y=-1$ 相交于点 $Q$。证明：以 $PQ$ 为直径的圆恒过 $y$ 轴上某定点。
\nopagebreak[4]
\begin{center}% ============================================================
% f2012_fujian_arts_q21_parabola.tex
% 2012 福建卷（文科）第 21 题 抛物线 E: x^2=2py (p>0)
% 等边三角形 OAB 的三个顶点均在抛物线上，A、B 关于 y 轴对称
% 来源 PDF：普通高考/2012/2012福建文.pdf，印面 732 第 1 栏中部。
% ============================================================
\begin{tikzpicture}[>=Stealth, line width=0.7pt, font=\small, x=1cm, y=1cm]
  \draw[->] (-3.0,0) -- (3.0,0) node[right] {$x$};
  \draw[->] (0,-1.0) -- (0,3.0) node[above] {$y$};

  \draw[thick, domain=-2.75:2.75, samples=100] plot (\x,{\x*\x/4});

  \coordinate (A) at (-2.40,1.44);
  \coordinate (B) at (2.40,1.44);
  \draw (A) -- (B);

  \fill (A) circle (1.2pt) node[left] {$A$};
  \fill (B) circle (1.2pt) node[right] {$B$};
  \fill (0,0) circle (1.2pt) node[below] {$O$};
\end{tikzpicture}
\end{center}
\nopagebreak[4]
\end{question}
\begin{question}
\gkti{22}已知函数 $f(x)=ax\sin x-\dfrac32$（$a\in\mathbf R$），且在 $\left[0,\dfrac\pi2\right]$ 上的最大值为 $\dfrac{\pi-3}{2}$。

（1） 求函数 $f(x)$ 的解析式；

（2） 判断函数 $f(x)$ 在 $(0,\pi)$ 内的零点个数，并加以证明。

\end{question}